\documentclass{article}
\usepackage[a4paper, margin=2.5cm]{geometry}
\usepackage{parskip}

\begin{document}

\begin{center}
    {\large \textbf{SURAT KUASA KHUSUS}} \\[0.5em]
    \rule{8cm}{0.4pt}
\end{center}

Yang bertanda tangan di bawah ini:
\begin{tabbing}
    \hspace{3cm} \= \hspace{1cm} \= \kill
    \textbf{Nama Lengkap} \> : \> Rahmat Hidayat \\
    \textbf{Pekerjaan} \> : \> Pegawai Swasta \\
    \textbf{Alamat} \> : \> Jl. Beringin, Kota Makassar \\
    \textbf{NIK} \> : \> 190627140674098
\end{tabbing}

Selanjutnya disebut sebagai \textbf{PEMBERI KUASA};

\vspace{1em}
\begin{center}
    \textbf{DENGAN INI MEMBERIKAN KUASA SEPENUHNYA KEPADA}
\end{center}
\vspace{0.5em}

\begin{tabbing}
    \hspace{3cm} \= \hspace{1cm} \= \kill
    \textbf{Nama Lengkap} \> : \> Liya, S.H. \\
    \textbf{Pekerjaan} \> : \> Advokat / Konsultan Hukum pada Kantor Hukum Justisia \\
    \textbf{Alamat} \> : \> Jl. Sukadamai, Kota Makassar \\
\end{tabbing}

Selanjutnya disebut sebagai \textbf{ PENERIMA KUASA };

\vspace{1em}
\begin{center}
    \textbf{------------------------- KHUSUS -------------------------}
\end{center}
\vspace{0.5em}

Untuk dan atas nama Pemberi Kuasa, bertindak sebagai Penggugat guna melakukan tindakan hukum mengajukan Gugatan Perbuatan Melawan Hukum terhadap \textbf{JOKO SAPUTRA} (selaku Tergugat) di Pengadilan Negeri Makassar berkaitan dengan iktikad buruk dan penyalahgunaan hak milik tanah dalam perbuatan pembangunan tembok batas setinggi 4 (empat) meter di lahan Tergugat yang mengakibatkan tertutupnya akses ventilasi dan jendela rumah Penggugat.

Untuk keperluan tersebut, Penerima Kuasa dikuasakan untuk:
\begin{enumerate}
    \item Menandatangani, mengajukan, dan mendaftarkan Surat Gugatan ke Kepaniteraan Pengadilan Negeri Makassar;
    \item Menghadiri persidangan-persidangan di Pengadilan Negeri Makassar, menghadap Hakim/Majelis Hakim, serta pejabat yang berwenang;
    \item Mengajukan dan menandatangani jawab-menjawab, replik, kesimpulan, serta menyerahkan alat-alat bukti surat maupun saksi-saksi;
    \item Meminta penetapan-penetapan, putusan, serta melaksanakan segala upaya hukum yang diperlukan demi kepentingan Pemberi Kuasa;
    \item Melakukan perdamaian jika dipandang perlu atas persetujuan Pemberi Kuasa, dan menandatangani akta perdamaian.
    \item Serta melakukan segala tindakan yang dianggap perlu oleh Penerima Kuasa untuk kepentingan Pemberi Kuasa sesuai dengan ketentuan Perundang-undangan yang berlaku, tanpa terkecuali.
\end{enumerate}

Demikian Surat Kuasa ini diberikan dengan hak substitusi, baik seluruhnya maupun sebagian, serta dengan tegas memberikan hak retensi menurut undang-undang.  

\vspace{1.5em}

\begin{flushright}
    Makassar, 17 September 2026 \\[2em]
\end{flushright}

\noindent
\begin{minipage}[t]{0.48\textwidth}
    Penerima Kuasa, \\[1em]
    % Kotak ilustrasi e-meterai darurat (tanpa asset gambar)
    \framebox{
        \begin{minipage}{3.2cm}
            \centering
            \tiny \textbf{E-MATERAI} \\
            \vspace{0.2em}
            \scriptsize Rp10.000 \\
            \tiny [ Sepuluh ribu rupiah ]
        \end{minipage}
    }\\[0.5em]
    \textbf{Liya, S.H.}
\end{minipage}
\hfill
\begin{minipage}[t]{0.48\textwidth}
    Pemberi Kuasa, \\[3.5em]
    \textbf{RAHMAT HIDAYAT}
\end{minipage}

\end{document}
